\section{国民年金}

\begin{fillblank}[自作 2026/11/1]
国民年金第１号被保険者に対して産前産後期間の保険料免除制度が施行され、
2019年4月分以後の国民年金の保険料月額（2004年度水準）は100円引き上がり（　　A　　）円となった。
\end{fillblank}

\begin{answerbox}
  \begin{enumerate}[A: ]
    \item 17,000
  \end{enumerate}
\end{answerbox}

\begin{explanationbox}
類題：年金1 2019 問題1(5)\textcircled{5}
\end{explanationbox}


\begin{truefalse}[年金2 2024 問題1(1)オ]
  下線部の正誤を判定し、誤っている場合は正しい内容に改めなさい。
  \\

厚生年金の保険料は、標準報酬月額と標準賞与額それぞれに対して一定率を賦課されるものとして
定められており、国民年金の保険料は、月額として定められている。2004(平成 16)年改正では、
これらを 2017(平成 29)年以降、\underline{厚生年金は 20％、国民年金は 27,500 円}
(2004(平成 16)年度価格)で固定することとされた。 (なお、厚生年金の保険料については、
これを労使折半することとされており、国民年金 の保険料については、2019(令和元)年度以降、
産前産後期間の保険料免除により固定額が これより引上げられている。)
\end{truefalse}

\begin{answerbox}
  \textbf{誤}

厚生年金は18.3％、国民年金は16,900円
\end{answerbox}

\begin{fillblank}[自作 2026/11/1]
  以下は国民年金の保険料免除等の老齢基礎年金額への反映割合を記載した表である。以下の空欄を全て埋めなさい。

  \begin{table}[htbp]
    \centering
    \begin{tabular}{|c|c|c|} \hline 
      免除等の種類 &  \begin{tabular}{c}平成21年4月以後の期間\\（国庫負担割合が2分の1の期間）\end{tabular}& \begin{tabular}{c}平成21年3月以前の期間\\（国庫負担割合が3分の1の期間）\end{tabular}\\ \hline
      全額免除 & （　　A　　） & 3分の1 \\ \hline
      4分の3免除 & （　　B　　） & 2分の1 \\ \hline
      半額免除 & （　　C　　） & 3分の2 \\ \hline
      4分の1免除 & （　　D　　） & 6分の5 \\ \hline
      学生納付特例 & \multicolumn{2}{|c|}{なし} \\ \hline
      保険料納付猶予 & \multicolumn{2}{|c|}{なし} \\ \hline
    \end{tabular}
  \end{table}
\end{fillblank}

\begin{answerbox}
  \begin{enumerate}[A: ]
    \item 2分の1
    \item 8分の5
    \item 4分の3
    \item 8分の7
  \end{enumerate}
\end{answerbox}

\begin{explanationbox}
類題：年金1 2018 問題1(4)\textcircled{3}
\end{explanationbox}

\begin{shortanswer}[年金1 2024 問題2(3)イ]
国民年金においては、本人の申し出により６０歳以上の者も加入することができる（任意加入
制度）。この任意加入の対象となる者はどのような者か簡潔に入力しなさい。
\end{shortanswer}

\begin{answerbox}
  \begin{enumerate}[(1)]
    \item 日本国内に住所を有する（または外国に居住するが日本国籍を持っている）、60 歳以上65 歳
未満の者（年金の資格期間を満たしていない場合は 70 歳未満の者まで）
  \item 老齢基礎年金の繰上げ支給を受けていない者
  \item 20 歳以上60 歳未満までの保険料の納付月数が480 月（40 年）未満の者
  \item 現在、厚生年金保険に加入していない者
  \end{enumerate}
\end{answerbox}

\begin{fillblank}[年金1 2023 問題1(5)イ(d)]
  国民年金制度の給付の種類に該当するものとして選択肢の中から当てはまるものの記号をす
べて挙げると（　　A　　）である。

\begin{table}[htbp]
  \centering
  \begin{tabular}{|llll|} \hline
    【選択肢】 & & & \\ 
    （ア）老齢基礎年金 & （イ）老齢厚生年金 & （ウ）付加年金 & （エ）寡婦年金 \\ 
    （オ）老齢手当金 & （カ）障害手当金 & （キ）遺族手当金 & （ク）死亡一時金 \\ \hline
  \end{tabular}
\end{table}
\end{fillblank}

\begin{answerbox}
  （ア）、（ウ）、（エ）、（ク）
\end{answerbox}

\newpage